% getting_started.tex -- CLI-first guide from installation to first sync

\chapter{Getting Started}

This chapter guides a new user through the Multi-repo Sync CLI. The documented
entry point is \cgscli{}.

The primary lifecycle is:

\begin{enumerate}
  \item \texttt{initialise(.cgs)} $\rightarrow$ clone all repositories
        $\rightarrow$ \texttt{READY}
  \item \texttt{initialise(.gts)} $\rightarrow$ load or restore a saved
        snapshot $\rightarrow$ \texttt{READY}
  \item \texttt{pull(.cgs|.gts)} $\rightarrow$ resynchronise an existing tree
  \item \texttt{checkout} / \texttt{add} / \texttt{commit} / \texttt{push} /
        \texttt{tag} $\rightarrow$ tree-wide Git operations
  \item \texttt{freeze} $\rightarrow$ write the next \texttt{.gts} snapshot
        and update the project-local \texttt{.lgr}
\end{enumerate}

% -----------------------------------------------------------------------
\section{Prerequisites}

\begin{itemize}
  \item Python~3.11 or later.
  \item Git~2.30 or later available on \texttt{PATH}.
  \item Working Git authentication for every remote used by the workspace.
\end{itemize}

% -----------------------------------------------------------------------
\section{Installation}

Install \cgspkg{} from source with the Pixi-managed environment:

\begin{lstlisting}[language=bash]
$ git clone https://github.com/flipoyo/ComplexGitSync.git
$ cd ComplexGitSync
$ pixi install
$ pixi run cgitsync --help
\end{lstlisting}

% -----------------------------------------------------------------------
\section{Create or Check a Project Spec}

A \texttt{.cgs} file is the local authoring spec for the repository tree. It
declares the project, each repository remote, each relative path, and optional
nested configuration discovery.

\begin{lstlisting}[language=bash]
$ pixi run cgitsync validate examples/complexgitsync.cgs
$ pixi run cgitsync print examples/complexgitsync.cgs
\end{lstlisting}

A valid declared tree prints a lifecycle line such as:

\begin{lstlisting}[language=bash]
DECLARED ready=false complete=true
\end{lstlisting}

% -----------------------------------------------------------------------
\section{Initialise the Workspace}

Use \texttt{initialise} as the canonical first command. From a \texttt{.cgs}
file it clones the full repository tree, validates it, writes the runtime
snapshot under \texttt{<workspace>/.cgitsync/state/}, and leaves the tree
\texttt{READY}.

\begin{lstlisting}[language=bash]
$ pixi run cgitsync initialise examples/complexgitsync.cgs
\end{lstlisting}

When \texttt{--output-path} is omitted, \cgscli{} uses the workspace-level
default \texttt{CGSPATH=../..}, then derives
\texttt{CGSHOME=CGSPATH/<project-name>}. The explicit equivalent is:

\begin{lstlisting}[language=bash]
$ pixi run cgitsync initialise examples/complexgitsync.cgs --output-path "$CGSPATH"
\end{lstlisting}

From a previously generated snapshot:

\begin{lstlisting}[language=bash]
$ pixi run cgitsync initialise .cgitsync/state/complexgitsync.gts
\end{lstlisting}

Successful \texttt{.cgs} initialisation prints the explicit workflow
\texttt{load->expand->validate->clone}, the per-repository
\texttt{git clone} action, the resolved root path, and a compact tree outline.
If a stale partial checkout or previous submodule link blocks the clone step,
\texttt{initialise} fails explicitly and prints
\texttt{Try clean-init method}.  Use \texttt{clean-init} to insert a purge
phase between validation and cloning:

\begin{lstlisting}[language=bash]
$ pixi run cgitsync clean-init examples/complexgitsync.cgs
\end{lstlisting}

The explicit workflow is
\texttt{load->expand->validate->purge->clone}.  The purge phase can also be
run on its own:

\begin{lstlisting}[language=bash]
$ pixi run cgitsync purge examples/complexgitsync.cgs
\end{lstlisting}

\texttt{purge} removes generated clone state from \texttt{CGSHOME}: immediate
child repository directories declared directly under the project root, project
\texttt{*.lgr} files, and the root \texttt{.gitmodules} file.

% -----------------------------------------------------------------------
\section{Inspect the Runtime State}

\begin{lstlisting}[language=bash]
$ pixi run cgitsync view_tree .cgitsync/state/complexgitsync.gts
$ pixi run cgitsync view_tree .cgitsync/state/complexgitsync.gts --depth 2 --collapse parent_2
$ pixi run cgitsync view_operation .cgitsync/state/complexgitsync.gts
\end{lstlisting}

If a command accepts an optional snapshot and none is provided, \cgscli{}
discovers the latest \texttt{.gts} from \texttt{CGSHOME/.cgitsync/state/}.
\texttt{CGSHOME} can be set explicitly; for \texttt{initialise(.cgs)} the
default is \texttt{CGSPATH=../..}, and later commands can also discover the
workspace by walking upward from the current directory.

% -----------------------------------------------------------------------
\section{Run the Git Cycle}

Once the tree is \texttt{READY}, the Multi-repo Sync commands operate on every
repository in deterministic order:

\begin{lstlisting}[language=bash]
$ pixi run cgitsync pull
$ pixi run cgitsync checkout feature/my-branch
$ pixi run cgitsync add
$ pixi run cgitsync commit "feat: update project CGS#1"
$ pixi run cgitsync push
$ pixi run cgitsync tag v1.2.3
$ pixi run cgitsync freeze release-2026.05
\end{lstlisting}

Pass \texttt{--gts <snapshot.gts>} when you want to pin a command to an
explicit snapshot, and pass \texttt{--dry-run} to \texttt{add},
\texttt{commit}, \texttt{push}, \texttt{tag}, or \texttt{freeze} to preview
the plan without mutating repositories.

% -----------------------------------------------------------------------
\section{Log Files}

Every CLI command writes a timestamped log file. On Linux the default location
is \texttt{~/.local/state/ComplexGitSync/logs/}; if
\texttt{XDG\_STATE\_HOME} is set, logs are written under that state directory
instead. The CLI prints the resolved path as \texttt{log\_file=<path>}.

% -----------------------------------------------------------------------
\section{Next Steps}

See Chapter~\ref{chap:user-guide} for the command reference, document formats,
preflight checks, and troubleshooting notes.
